% HP-D10-CH28-C stable-unit-id: D10-FREMLIN-V2-S285
\frfilename{mt285.tex}
\versiondate{18.9.14}
\copyrightdate{1994}

\def\chaptername{Analisis Fourier}
\def\sectionname{Fungsi karakteristik}

\newsection{285}

Sekarang kita sampai pada salah satu penerapan transformasi Fourier yang paling
efektif, yakni penggunaan `fungsi karakteristik' untuk menganalisis distribusi
probabilitas. Ternyata transformasi Fourier suatu distribusi probabilitas bukan
hanya menentukan distribusi itu (285M), tetapi juga memudahkan perhitungan
banyak hal yang ingin kita ketahui tentang distribusi tersebut dari
transformasinya\cmmnt{ (285G, 285Xi)}. Bahkan yang lebih mencolok, konvergensi
titik demi titik transformasi Fourier bersesuaian (untuk barisan) dengan
konvergensi dalam topologi samar pada ruang distribusi. Karena itu, fungsi
karakteristik menyediakan metode baru yang sangat ampuh untuk membuktikan hasil
seperti Teorema Limit Pusat dan teorema Poisson (285Q).

Karena penerapan gagasan-gagasan di sini terutama termasuk dalam teori
probabilitas, saya kembali menggunakan istilah para ahli probabilitas, seperti
dalam Bab 27. Meskipun demikian, pada banyak bagian kita perlu berbicara tentang
integral, dan sering kali ada lebih dari satu ukuran yang terlibat. Karena itu,
saya tegaskan bahwa integral $\int f(x)dx$ selalu diambil terhadap ukuran
Lebesgue (biasanya, tetapi tidak selalu, satu dimensi), seperti pada bagian lain
bab ini, sedangkan integral terhadap ukuran lain ditulis dalam bentuk
$\int fd\nu$ atau $\int f(x)\nu(dx)$.

\leader{285A}{Definisi (a)} Misalkan $\nu$ suatu ukuran probabilitas Radon pada
$\BbbR^r$\cmmnt{ (256A)}. Maka {\bf fungsi karakteristik}
dari $\nu$ ialah fungsi $\phi_{\nu}:\BbbR^r\to\Bbb C$ yang diberikan oleh rumus

\Centerline{$\phi_{\nu}(y)=\int e^{iy\dotproduct x}\nu(dx)$}

\noindent untuk setiap $y\in\BbbR^r$, dengan menulis
$y\dotproduct x=\eta_1\xi_1+\ldots+\eta_r\xi_r$ jika
$y=(\eta_1,\ldots,\eta_r)$ dan $x=(\xi_1,\ldots,\xi_r)$.

\header{285Ab}{\bf (b)} Misalkan $X_1,\ldots,X_r$ variabel-variabel acak
bernilai real pada ruang probabilitas yang sama. {\bf Fungsi karakteristik}
dari $\pmb{X}=(X_1,\ldots,X_r)$ ialah fungsi karakteristik
$\phi_{\pmb{X}}=\phi_{\nu_{\pmb{X}}}$ dari distribusi probabilitas bersamanya
$\nu_{\pmb{X}}$, sebagaimana didefinisikan dalam 271C.

\cmmnt{
\leader{285B}{Catatan (a)} Karena salah satu kebetulan sejarah yang lazim,
definisi `fungsi karakteristik' dan `transformasi Fourier' berkembang secara
terpisah. Dalam 283Ba saya mencatat bahwa definisi transformasi Fourier belum
baku, dan bahwa rumus

\Centerline{$\varhatf(y)=\int_{-\infty}^{\infty}e^{iyx}f(x)dx$,}

\Centerline{$\varcheckf(y)
=\Bover1{2\pi}\int_{-\infty}^{\infty}e^{-iyx}f(x)dx$}

\noindent akan lebih mudah dipadankan dengan definisi fungsi karakteristik.
Saya tidak mengambil pilihan ini karena ingin menonjolkan simetri antara
transformasi Fourier dan transformasi Fourier invers, serta korespondensi antara
transformasi Fourier dan deret Fourier. Keuntungan utama menyamakan kedua
definisi, yaitu menyamakan konstanta dalam hasil seperti 283F dan 285Xk, akan
diimbangi oleh keharusan mengingat konstanta yang berbeda untuk $\varhatf$ dan
$\varcheckf$ dalam hasil seperti 283M.

\header{285Bb}{\bf (b)} Alasan lain untuk tidak memaksakan agar rumus-rumus
dalam bagian ini langsung bersesuaian dengan rumus-rumus dalam \S\S283-284
ialah bahwa kasus berdimensi $r$ sangat penting dalam beberapa penerapan fungsi
karakteristik yang utama, sehingga layak diperkenalkan sejak awal. Akibatnya,
rumus-rumus dalam bagian ini tetap akan mempunyai unsur baru dibandingkan dengan
rumus-rumus yang telah muncul sebelumnya.
}

\leader{285C}{}\cmmnt{ Tentu ada cara langsung untuk menggambarkan fungsi
karakteristik suatu keluarga $(X_1,\ldots,X_r)$ variabel acak, sebagai berikut.

\medskip

\noindent}{\bf Proposisi} Misalkan $X_1,\ldots,X_r$ variabel-variabel acak
bernilai real pada ruang probabilitas yang sama, dan $\nu_{\pmb{X}}$ distribusi
bersamanya. Maka fungsi karakteristiknya $\phi_{\nu_{\pmb{X}}}$ diberikan oleh

\Centerline{$\phi_{\nu_{\pmb{X}}}(y)=\Expn(e^{iy\dotproduct \pmb{X}})
=\Expn(e^{i\eta_1X_1}e^{i\eta_2X_2}\ldots e^{i\eta_rX_r})$}

\noindent untuk setiap $y=(\eta_1,\ldots,\eta_r)\in\BbbR^r$.

\proof{ Terapkan 271E pada fungsi $h_1$, $h_2:\Bbb R^r\to\Bbb R$
yang didefinisikan oleh

\Centerline{$h_1(x)=\cos(y\dotproduct x)$,\quad
$h_2(x)=\sin(y\dotproduct x)$,}

\noindent untuk memperoleh

$$\eqalign{\phi_{\nu_{\pmb{X}}}(y)
&=\int h_1(x)\nu_{\pmb{X}}(dx) + i\int h_2(x)\nu_{\pmb{X}}(dx)\cr
&=\Expn(h_1(\pmb{X})) + i\Expn(h_2(\pmb{X}))
=\Expn(e^{iy\dotproduct \pmb{X}}).\cr}$$
}%end of proof of 285C

\vleader{72pt}{285D}{}\cmmnt{ Saya perlu menjelaskan korespondensi antara
transformasi Fourier, sebagaimana didefinisikan dalam 283A, dan fungsi
karakteristik.

\medskip

\noindent}{\bf Proposisi} Misalkan $\nu$ suatu ukuran probabilitas Radon pada
$\Bbb R$. Tuliskan

\Centerline{$\varhat{\nu}(y)
=\Bover1{\sqrt{2\pi}}\int_{-\infty}^{\infty}
e^{-iyx}\nu(dx)$}

\noindent untuk setiap $y\in\Bbb R$, dan tuliskan $\phi_{\nu}$ untuk fungsi
karakteristik dari $\nu$.


(a) $\varhat{\nu}(y)=\Bover1{\sqrt{2\pi}}\phi_{\nu}(-y)$ untuk setiap
$y\in\Bbb R$.

(b) Untuk setiap fungsi Lebesgue terintegralkan bernilai kompleks $h$ pada
$\Bbb R$,

\Centerline{$\int_{-\infty}^{\infty}\varhat{\nu}(y)h(y)dy
=\int_{-\infty}^{\infty}\varhat{h}(x)\nu(dx)$.}

(c) Untuk setiap fungsi uji yang menurun cepat $h$ pada $\Bbb R$\cmmnt{ (lihat
\S284)},

\Centerline{$\int_{-\infty}^{\infty}h(x)\nu(dx)
=\int_{-\infty}^{\infty}\varcheck{h}(y)\varhat{\nu}(y)dy$.}

(d) Jika $\nu$ merupakan ukuran integral tak tentu terhadap ukuran Lebesgue,
dengan turunan Radon--Nikod\'ym $f$, maka $\varhat{\nu}$ merupakan transformasi
Fourier dari $f$.

\proof{{\bf (a)} Ini langsung mengikuti dari definisi
$\phi_{\nu}$ dan $\varhat{\nu}$.

\medskip

{\bf (b)} Karena

\Centerline{$\int_{-\infty}^{\infty}\int_{-\infty}^{\infty}
|h(y)|\nu(dx)dy=\int_{-\infty}^{\infty}|h(y)|dy<\infty$,}

\noindent kita dapat menukar urutan integrasi dan memperoleh

$$\eqalign{\int_{-\infty}^{\infty}\varhat{\nu}(y)h(y)dy
&=\Bover1{\sqrt{2\pi}}\int_{-\infty}^{\infty}\int_{-\infty}^{\infty}
e^{-iyx}h(y)\nu(dx)dy\cr
&=\Bover1{\sqrt{2\pi}}\int_{-\infty}^{\infty}\int_{-\infty}^{\infty}
e^{-iyx}h(y)dy\,\nu(dx)
=\int_{-\infty}^{\infty}\varhat{h}(x)\nu(dx).\cr}$$

\medskip

{\bf (c)} Ini langsung mengikuti dari (b), karena $\varcheck{h}$ terintegralkan
dan $\varcheck{h}\varhat{\phantom{h}}=h$ (284C).

\medskip

{\bf (d)} Intinya adalah bahwa

\Centerline{$\int h\,d\nu=\int h(x)f(x)dx$}

\noindent untuk setiap fungsi Borel terukur terbatas $h:\Bbb R\to\Bbb R$
(235K), dan karenanya juga untuk fungsi
$x\mapsto e^{-iyx}:\Bbb R\to\Bbb C$. Dengan demikian

\Centerline{$\varhat{\nu}(y)
=\Bover1{\sqrt{2\pi}}\int_{-\infty}^{\infty}e^{-iyx}\nu(dx)
=\Bover1{\sqrt{2\pi}}\int_{-\infty}^{\infty}e^{-iyx}f(x)dx
=\varhatf(y)$}

\noindent untuk setiap $y$.
}% end of proof of 285D

\leader{285E}{Lema} Misalkan $X$ suatu variabel acak normal dengan ekspektasi
$a$ dan varians $\sigma^2$, dengan $\sigma>0$. Maka fungsi karakteristik $X$
diberikan oleh rumus

\Centerline{$\phi(y)=e^{iya}e^{-\sigma^2y^2/2}$.}

\proof{ Ini hanyalah 283N dengan konstanta yang disesuaikan.

$$\eqalignno{\phi(y)
&=\Expn(e^{iyX})
=\Bover1{\sigma\sqrt{2\pi}}\int_{-\infty}^{\infty}
  e^{iyx}e^{-(x-a)^2/2\sigma^2}dx\cr
\noalign{\noindent (dengan menggunakan fungsi kepadatan $X$ yang diberikan
dalam 274Ad dan menerapkan 271Ic)}
&=\Bover1{\sqrt{2\pi}}\int_{-\infty}^{\infty}
  e^{iy(\sigma t+a)}e^{-t^2/2}dt\cr
\noalign{\noindent (dengan menyubstitusikan $x=\sigma t+a$)}
&=e^{iya}\sqrt{2\pi}\varhat{\psi}_1(-y\sigma)\cr
\noalign{\noindent (dengan menetapkan
$\psi_1(x)=\bover1{\sqrt{2\pi}}e^{-x^2/2}$, seperti dalam 283N)}
&=e^{iya}e^{-\sigma^2y^2/2}.\cr}$$
}%end of proof of 285E

\leader{285F}{}\cmmnt{ Sekarang saya berikan hasil yang bersesuaian dengan
sebagian 283C, dengan penyempurnaan tambahan mengenai variabel acak bebas
(285I).

\medskip

\noindent}{\bf Proposisi} Misalkan $\nu$ suatu ukuran probabilitas Radon pada
$\BbbR^r$, dan $\phi$ fungsi karakteristiknya.

(a) $\phi(0)=1$.

(b) $\phi:\BbbR^r\to\Bbb C$ kontinu seragam.

(c) $\phi(-y)=\overline{\phi(y)}$, $|\phi(y)|\le 1$ untuk setiap $y\in\Bbb
R^r$.

(d) Jika $r=1$ dan $\int|x|\nu(dx)<\infty$, maka $\phi'(y)$ ada dan sama dengan $i\int xe^{ixy}\nu(dx)$ untuk setiap $y\in\Bbb R$.

(e) Jika $r=1$ dan $\int x^2\nu(dx)<\infty$, maka $\phi''(y)$ ada dan sama dengan $-\int x^2e^{ixy}\nu(dx)$ untuk setiap $y\in\Bbb R$.

\proof{{\bf (a)} $\phi(0)=\int\chi\BbbR^r\,d\nu=\nu(\BbbR^r)=1$.

\medskip

{\bf (b)} Misalkan $\epsilon>0$. Pilih $M>0$ sedemikian sehingga

\Centerline{$\nu\{x:\|x\|\ge M\}\le\epsilon$,}

\noindent dengan menulis $\|x\|=\sqrt{x\dotproduct x}$ seperti biasa. Pilih
$\delta>0$ sedemikian sehingga $|e^{ia}-1|\le\epsilon$ bila
$|a|\le\delta$. Sekarang andaikan $y$, $y'\in\BbbR^r$ memenuhi
$\|y-y'\|\le\delta/M$. Maka, bila $\|x\|\le M$,

\Centerline{$|e^{iy\dotproduct x}-e^{iy'\dotproduct x}|
=|e^{iy'\dotproduct x}||e^{i(y-y')\dotproduct x}-1|
=|e^{i(y-y')\dotproduct x}-1|
\le\epsilon$}

\noindent karena

\Centerline{$|(y-y')\dotproduct x|\le\|y-y'\|\|x\|\le\delta$.}

\noindent Oleh karena itu, menulis $B$ untuk $\{x:\|x\|\le M\}$,

$$\eqalign{|\phi(y)-\phi(y')|
&\le\int_B|e^{iy\dotproduct x}-e^{iy'\dotproduct x}|\nu(dx)
\cr&\qquad\qquad
  +\int_{\BbbR^r\setminus B}|e^{iy\dotproduct x}|\nu(dx)
  +\int_{\BbbR^r\setminus B}|e^{iy'\dotproduct x}|\nu(dx)\cr
&\le\epsilon+\epsilon+\epsilon
=3\epsilon.\cr}$$

\noindent Karena $\epsilon$ sembarang, $\phi$ kontinu seragam.


\medskip

{\bf (c)} Ini langsung;

\Centerline{$\phi(-y)=\int e^{-iy\dotproduct x}\nu(dx)
=\overline{\intop e^{iy\dotproduct x}\nu(dx)}=\overline{\phi(y)}$,}

\Centerline{$|\phi(y)|=|\int e^{iy\dotproduct x}\nu(dx)|
\le\int|e^{iy\dotproduct x}|\nu(dx)
=1$.}

\medskip

{\bf (d)} Intinya adalah bahwa $|\pd{}{y}e^{iyx}|=|x|$ untuk setiap $x$,
$y\in\Bbb R$. Jadi, menurut 123D (diterapkan, tepatnya, pada bagian real dan
imajiner fungsi),

\Centerline{$\phi'(y)
=\Bover{d}{dy}\int e^{iyx}\nu(dx)
=\int\Pd{}{y}e^{iyx}\nu(dx)
=\int ixe^{iyx}\nu(dx)$.}

\medskip

{\bf (e)} Karena kini $|\pd{}{y}xe^{iyx}|=x^2$ untuk setiap $x$, $y$, kita
dapat mengulangi argumen untuk memperoleh

\Centerline{$\phi''(y)
=i\Bover{d}{dy}
     \int xe^{iyx}\nu(dx)
=i\int
     \Pd{}{y}xe^{iyx}\nu(dx)
=-\int
     x^2e^{iyx}\nu(dx)$.}
}%end of proof of 285F

\vleader{48pt}{285G}{Akibat} (a) Misalkan $X$ variabel acak bernilai real dengan
ekspektasi berhingga, dan $\phi$ fungsi karakteristiknya.
$\phi'(0)=i\Expn(X)$.

(b) Misalkan $X$ variabel acak bernilai real dengan varians berhingga, dan
$\phi$ fungsi karakteristiknya.
$\phi''(0)=-\Expn(X^2)$.

\proof{ Kita hanya perlu mencocokkan $X$ dengan distribusinya $\nu$, lalu
mengamati bahwa

\Centerline{`$X$ mempunyai ekspektasi berhingga'}

\noindent bersesuaian dengan

\Centerline{`$\int |x|\nu(dx)=\Expn(|X|)<\infty$',}

\noindent jadi bahwa

\Centerline{$\phi'(0)=i\int x\,\nu(dx)=i\Expn(X)$,}

\noindent dan bahwa

\Centerline{`$X$ mempunyai varians berhingga'}

\noindent bersesuaian dengan

\Centerline{`$\int x^2\nu(dx)=\Expn(X^2)<\infty$',}

\noindent jadi bahwa

\Centerline{$\phi''(0)=-\int x^2\,\nu(dx)=-\Expn(X^2)$,}

\noindent sebagaimana dalam 271E.
}%end of proof of 285G

\cmmnt{
\leader{285H}{Catatan} Perhatikan bahwa tidak ada hasil yang bersesuaian dengan
283Cg (`$\lim_{|y|\to\infty}\varhatf(y)
\ifdim\pagewidth=390pt\penalty-100\fi
=0$'). Jika $\nu$ adalah ukuran Dirac pada $\Bbb R$ yang terpusat di $0$,
yakni distribusi suatu variabel acak yang hampir di mana-mana bernilai nol,
maka $\phi(y)=1$ untuk setiap $y$.
}

\leader{285I}{Proposisi} Misalkan $X_1,\ldots,X_n$ variabel-variabel acak
bernilai real yang bebas, dengan fungsi karakteristik
$\phi_{1},\ldots,\phi_{n}$. Misalkan $\phi$ fungsi karakteristik dari jumlahnya
$X=X_1+\ldots+X_n$. Maka

\Centerline{$\phi(y)=\prod_{j=1}^n\phi_j(y)$}

\noindent untuk setiap $y\in\Bbb R$.

\proof{ Misalkan $y\in\Bbb R$. Menurut 272E, variabel-variabel

\Centerline{$Y_j=e^{iyX_j}$}

\noindent saling bebas, sehingga menurut 272R

\Centerline{$\phi(y)=\Expn(e^{iyX})=\Expn(e^{iy(X_1+\ldots+X_n)})
=\Expn(\prod_{j=1}^nY_j)=\prod_{j=1}^n\Expn(Y_j)
=\prod_{j=1}^n\phi_j(y)$,}

\noindent sesuai kebutuhan.
}%end of proof of 285I

\cmmnt{
\medskip

\noindent{\bf Komentar} Lihat juga 285R di bawah.
}

\leader{285J}{}\cmmnt{ Ada teorema inversi untuk fungsi karakteristik yang
bersesuaian dengan 283F. Saya menyajikannya dalam 285Xk, dengan versi
berdimensi $r$ dalam 285Yb. Namun, hasil ini tampaknya tidak seberguna rangkaian
hasil berikut.

\medskip

\noindent}{\bf Lema} Misalkan $\nu$ suatu ukuran probabilitas Radon pada
$\Bbb R^r$, dan $\phi$ fungsi karakteristiknya. Maka untuk $1\le j\le r$ dan
$a>0$,

\Centerline{$\nu\{x:|\xi_j|\ge a\}
\le 7a\int_0^{1/a}(1-\Real\phi(te_j))dt$,}

\noindent dengan $e_j\in\BbbR^r$ vektor satuan ke-$j$.

\wheader{285J}{0}{0}{0}{84pt}

\proof{ Kita mempunyai

$$\eqalignno{7a\int_0^{1/a}(1-\Real\phi(te_j))dt
&=7a\int_0^{1/a}
   \bigl(1-\Real\int_{\Bbb R^r}e^{it\xi_j}\nu(dx)\bigr)dt\cr
&=7a\int_0^{1/a}\int_{\Bbb R^r}1-\cos(t\xi_j)\nu(dx)dt\cr
&=7a\int_{\Bbb R^r}\int_0^{1/a}1-\cos(t\xi_j)dt\,\nu(dx)\cr
\displaycause{karena $(x,t)\mapsto 1-\cos(t\xi_j)$ terbatas dan
$\nu\BbbR^r\cdot\Bover1a$ berhingga}
&=7a\int_{\Bbb R^r}
  \bigl(\bover1a-\bover1{\xi_j}\sin\bover{\xi_j}a\bigr)\nu(dx)\cr
&\ge7a\int_{|\xi_j|\ge a}
\bigl(\bover1a-\bover1{\xi_j}\sin\bover{\xi_j}a\bigr)\nu(dx)\cr
\noalign{\noindent (karena $\Bover1{\xi}\sin\Bover{\xi}{a}\le\Bover1a$
untuk setiap $\xi\ne 0$)}
&\ge\nu\{x:|\xi_j|\ge a\},\cr}$$

\noindent karena

\Centerline{$\Bover{\sin\eta}{\eta}\le\Bover{\sin 1}{1}\le\Bover67$ jika
$\eta\ge 1$,}

\noindent sehingga

\Centerline{$a(\Bover1a-\Bover1{\xi_j}\sin\Bover{\xi_j}a)\ge\Bover17$}

\noindent jika $|\xi_j|\ge a$.
}%end of proof of 285J

\cmmnt{
\leader{285K}{Fungsi karakteristik dan topologi samar}
Kini saatnya kembali pada gagasan yang disebut secara singkat dalam 274L.
Tetapkan $r\ge 1$ dan misalkan $P$ himpunan semua ukuran probabilitas Radon
pada $\BbbR^r$. Untuk setiap fungsi kontinu terbatas
$h:\BbbR^r\to\Bbb R$, definisikan $\rho_h:P\times P\to\Bbb R$ dengan

\Centerline{$\rho_h(\nu,\nuprime)=|\int h\,d\nu-\int h\,d\nuprime|$}

\noindent untuk $\nu$, $\nuprime\in P$. Maka topologi samar pada $P$ ialah
topologi yang dibangkitkan oleh pseudometrik-pseudometrik $\rho_h$ (274Ld).
}

\leader{285L}{Teorema} Misalkan $\nu$, $\sequencen{\nu_n}$ ukuran-ukuran
probabilitas Radon pada $\BbbR^r$, dengan fungsi karakteristik $\phi$,
$\sequencen{\phi_n}$. Maka pernyataan-pernyataan berikut ekuivalen:

(i) $\nu=\lim_{n\to\infty}\nu_n$ dalam topologi samar;

(ii) $\int h\,d\nu=\lim_{n\to\infty}\int h\,d\nu_n$ untuk setiap fungsi
kontinu terbatas $h:\BbbR^r\to\Bbb R$;

(iii) $\lim_{n\to\infty}\phi_n(y)=\phi(y)$ untuk setiap $y\in\BbbR^r$.

\proof{{\bf (a)} Ekuivalensi (i) dan (ii) pada dasarnya merupakan definisi
topologi samar; kita mempunyai

$$\eqalignno{\lim_{n\to\infty}
&\nu_n=\nu\text{ dalam topologi samar}\cr
&\iff\,\lim_{n\to\infty}\rho_h(\nu_n,\nu)=0
\text{ untuk setiap fungsi kontinu terbatas }h\cr
\noalign{\noindent (2A3Mc)}
&\iff\,\lim_{n\to\infty}|\int h\,d\nu_n-\int h\,d\nu|=0
   \text{ untuk setiap fungsi kontinu terbatas }h.\cr}$$

\medskip

{\bf (b)} Selanjutnya, (ii) jelas mengakibatkan (iii), karena

\Centerline{$\Real\phi(y)
=\int h_y\,d\nu
=\lim_{n\to\infty}\int h_y\,d\nu_n
=\lim_{n\to\infty}\Real\phi_n(y)$,}

\noindent dengan menetapkan $h_y(x)=\cos(x\dotproduct y)$ untuk setiap $x$,
dan demikian pula

\Centerline{$\Imag\phi(y)=\lim_{n\to\infty}\Imag\phi_n(y)$}

\noindent untuk setiap $y\in\BbbR^r$.

\medskip

{\bf (c)} Jadi tinggal dibuktikan bahwa (iii)$\Rightarrow$(ii). Saya mulai
dengan menunjukkan bahwa, untuk setiap $\epsilon>0$, ada himpunan tertutup
terbatas $K$ sedemikian sehingga

\Centerline{$\nu_n(\BbbR^r\setminus K)\le\epsilon$ untuk setiap
$n\in\Bbb N$.}

\noindent\Prf\ Kita tahu bahwa $\phi(0)=1$ dan $\phi$ kontinu di
$0$ (285Fb). Pilih $a>0$ cukup besar sehingga bila $j\le r$ dan
$|t|\le 1/a$, berlaku

\Centerline{$1-\Real\phi(te_j)\le\Bover{\epsilon}{14r}$,}

\noindent dengan menulis $e_j$ untuk vektor satuan ke-$j$, seperti dalam
285J. Maka

\Centerline{$7a\int_0^{1/a}
(1-\Real\phi(te_j))dt\le\Bover{\epsilon}{2r}$}

\noindent untuk setiap $j\le r$. Menurut Teorema Konvergensi Terdominasi
Lebesgue (karena fungsi-fungsi $t\mapsto 1-\Real\phi_n(te_j)$ tentu terbatas
secara seragam pada $[0,\bover1a]$), ada $n_0\in\Bbb N$ sedemikian sehingga

\Centerline{$7a\int_0^{1/a}
(1-\Real\phi_n(te_j))dt\le\Bover{\epsilon}{r}$}

\noindent untuk setiap $j\le r$ dan $n\ge n_0$. Namun, 285J sekarang memberi

\Centerline{$\nu_n\{x:|\xi_j|\ge a\}\le\Bover{\epsilon}{r}$}

\noindent untuk setiap $j\le r$, $n\ge n_0$. Di sisi lain, tentu ada $b\ge a$
sedemikian sehingga

\Centerline{$\nu_n\{x:|\xi_j|\ge b\}\le\Bover{\epsilon}{r}$}

\noindent untuk setiap $j\le r$, $n<n_0$. Jadi, dengan menetapkan
$K=\{x:|\xi_j|\le b\text{ untuk setiap }j\le r\}$,

\Centerline{$\nu_n(\BbbR^r\setminus K)\le\epsilon$}

\noindent untuk setiap $n\in\Bbb N$, sesuai kebutuhan.  \Qed

\medskip

{\bf (d)} Sekarang ambil sembarang fungsi kontinu terbatas
$h:\BbbR^r\to\Bbb R$ dan $\epsilon>0$. Tetapkan
$M=1+\sup_{x\in\BbbR^r}|h(x)|$, dan pilih himpunan tertutup terbatas $K$
sedemikian sehingga

\Centerline{$\nu_n(\BbbR^r\setminus K)\le\Bover{\epsilon}M$ untuk setiap
$n\in\Bbb N$, \quad$\nu(\BbbR^r\setminus K)\le\Bover{\epsilon}M$,}

\noindent dengan menggunakan (b) di atas. Menurut teorema Stone--Weierstrass
(281K), ada $y_0,\ldots,y_m\in\Bbb Q^r$ dan
$c_0,\ldots,c_m\in\Bbb C$ sedemikian sehingga

\Centerline{$|h(x)-g(x)|\le\epsilon$ untuk setiap $x\in K$,}

\Centerline{$|g(x)|\le M$ untuk setiap $x\in\BbbR^r$,}

\noindent dengan menulis $g(x)=\sum_{k=0}^mc_ke^{iy_k\dotproduct x}$ untuk
$x\in\BbbR^r$. Sekarang

\Centerline{$\lim_{n\to\infty}\int g\,d\nu_n
=\lim_{n\to\infty}\sum_{k=0}^mc_k\phi_n(y_k)
=\sum_{k=0}^mc_k\phi(y_k)
=\int g\,d\nu$.}

\noindent Di sisi lain, untuk setiap $n\in\Bbb N$,

\Centerline{$|\int g\,d\nu_n-\int h\,d\nu_n|
\le\int_K|g-h|d\nu_n+2M\nu_n(\BbbR^r\setminus K)
\le 3\epsilon$,}

\noindent dan dengan cara yang sama
$|\int g\,d\nu-\int h\,d\nu|\le 3\epsilon$. Oleh karena itu,

\Centerline{$\limsup_{n\to\infty}|\int h\,d\nu_n-\int h\,d\nu|\le
6\epsilon$.}

\noindent Karena $\epsilon$ sembarang,

\Centerline{$\lim_{n\to\infty}\int h\,d\nu_n=\int h\,d\nu$,}

\noindent dan (ii) benar.
}%end of proof of 285L


\vleader{72pt}{285M}{Akibat} (a) Misalkan $\nu$, $\nuprime$ dua ukuran
probabilitas Radon pada $\BbbR^r$ dengan fungsi karakteristik yang sama. Maka
keduanya sama.

(b) Misalkan $(X_1,\ldots,X_r)$ dan $(Y_1,\ldots,Y_r)$ dua keluarga variabel
acak bernilai real. Jika

\Centerline{$\Expn(e^{i\eta_1X_1+\ldots+i\eta_rX_r})
=\Expn(e^{i\eta_1Y_1+\ldots+i\eta_rY_r})$}

\noindent untuk semua $\eta_1,\ldots,\eta_r\in\Bbb R$, maka
$(X_1,\ldots,X_r)$ memiliki distribusi bersama yang sama dengan
$(Y_1,\ldots,Y_r)$.

\proof{{\bf (a)} Terapkan 285L dengan $\nu_n=\nuprime$ untuk setiap $n$.
Kita memperoleh $\int h\,d\nuprime=\int h\,d\nu$ untuk setiap fungsi kontinu
terbatas $h:\BbbR^r\to\Bbb R$. Menurut 256D(iv), $\nu=\nuprime$.

\medskip

{\bf (b)} Terapkan (a) dengan $\nu$, $\nuprime$ sebagai kedua distribusi
bersama tersebut.
}%end of proof of 285M

\cmmnt{
\leader{285N}{Catatan} Penerapan terpenting teorema ini barangkali adalah
pembuktian baku Teorema Limit Pusat. Saya menguraikan gagasannya dalam 285Xq
dan 285Yl-285Yo; rinciannya dapat ditemukan dalam kebanyakan buku teori
probabilitas yang serius. Dua buku di rak saya ialah {\smc Shiryayev 84},
\S III.4, dan {\smc Feller 66}, \S XV.6. Namun, untuk memperoleh kekuatan penuh
versi Lindeberg dari Teorema Limit Pusat, kita harus bekerja cukup keras. Karena
itu, sebagai gantinya saya akan menggambarkan metode tersebut melalui suatu
versi teorema Poisson (285Q). Saya mulai dengan dua lema yang sangat sering
digunakan dalam hasil semacam ini.
}

\leader{285O}{Lema} Misalkan
$c_0,\ldots,c_n,d_0,\ldots,d_n$ bilangan-bilangan kompleks dengan modulus
paling besar $1$. Maka

\Centerline{$|\prod_{k=0}^nc_k-\prod_{k=0}^nd_k|
\le\sum_{k=0}^n|c_k-d_k|$.}

\proof{ Gunakan induksi pada $n$. Kasus $n=0$ langsung. Untuk kasus $n=1$,
kita mempunyai

$$\eqalign{|c_0c_1-d_0d_1|
&=|c_0(c_1-d_1)+(c_0-d_0)d_1|\cr
&\le|c_0||c_1-d_1|+|c_0-d_0||d_1|
\le|c_1-d_1|+|c_0-d_0|,\cr}$$

\noindent sebagaimana diperlukan. Untuk langkah induksi ke $n+1$, kita
mempunyai

$$\eqalignno{|\prod_{k=0}^{n+1}c_k-\prod_{k=0}^{n+1}d_k|
&\le|\prod_{k=0}^nc_k-\prod_{k=0}^nd_k|
+|c_{n+1}-d_{n+1}|\cr
\noalign{\noindent (menurut kasus yang baru dibuktikan, karena
$c_{n+1}$, $d_{n+1}$, $\prod_{k=0}^nc_k$, dan $\prod_{k=0}^nd_k$ semuanya
bermodulus paling besar $1$)}
&\le\sum_{k=0}^n|c_k-d_k|+|c_{n+1}-d_{n+1}|\cr
\noalign{\noindent (menurut hipotesis induksi)}
&=\sum_{k=0}^{n+1}|c_k-d_k|,\cr}$$

\noindent sehingga induksi berlanjut.
}%end of proof of 285O

\leader{285P}{Lema} Misalkan $M\ge 0$ dan $\epsilon>0$. Maka ada
$\eta>0$ dan $y_0,\ldots,y_n\in\Bbb R$ sedemikian sehingga, bila $X$, $Z$
dua variabel acak bernilai real dengan $\Expn(|X|)\le M$,
$\Expn(|Z|)\le M$ dan
$|\phi_X(y_j)-\phi_Z(y_j)|\le\eta$ untuk setiap $j\le n$, maka
$F_X(a)\le F_Z(a+\epsilon)+\epsilon$ untuk setiap $a\in\Bbb R$, dengan
$\phi_X$ menyatakan fungsi karakteristik $X$ dan $F_X$ menyatakan fungsi
distribusi $X$.

\proof{ Kasus $M=0$ langsung, sebab ketika itu $X$ dan $Z$ keduanya nol hampir
di mana-mana. Jadi selanjutnya saya andaikan $M>0$. Tetapkan
$\delta=\Bover{\epsilon}7>0$, $b=\Bover{M}{\delta}$.

\medskip

{\bf (a)} Definisikan $h_0:\Bbb R\to[0,1]$ dengan
$h_0(x)=\med(0,1-\Bover{x}{\delta},1)$ untuk $x\in\Bbb R$. Maka $h_0$
kontinu. Misalkan $m=\lfloor\bover{b}{\delta}\rfloor$ merupakan bagian bulat
dari $\bover{b}{\delta}$, dan untuk $-m\le k\le m+1$ tetapkan
$h_k(x)=h_0(x-k\delta)$.

Menurut teorema Stone--Weierstrass (sekali lagi 281K), ada
$y_0,\ldots,y_n\in\Bbb R$ dan $c_0,\ldots,c_n\in\Bbb C$ sedemikian sehingga,
dengan menulis $g_0(x)=\sum_{j=0}^nc_je^{iy_jx}$,

\Centerline{$|h_0(x)-g_0(x)|\le\delta$
  untuk setiap $x\in[-b-(m+1)\delta,b+m\delta]$,}

\Centerline{$|g_0(x)|\le 1$ untuk setiap $x\in\Bbb R$.}

\noindent Untuk $-m\le k\le m+1$, tetapkan

\Centerline{$g_k(x)=g_0(x-k\delta)
=\sum_{j=0}^nc_je^{-iy_jk\delta}e^{iy_jx}$.}

\noindent Tetapkan $\eta=\delta/(1+\sum_{j=0}^n|c_j|)>0$.

\medskip

{\bf (b)} Sekarang andaikan $X$, $Z$ variabel-variabel acak yang memenuhi
$\Expn(|X|)\le M$, $\Expn(|Z|)\le M$ dan
$|\phi_X(y_j)-\phi_Z(y_j)|\le\eta$ untuk setiap $j\le n$. Maka untuk setiap
$k$ kita mempunyai

\Centerline{$\Expn(g_k(X))
=\Expn(\sum_{j=0}^nc_je^{-iy_jk\delta}e^{iy_jX})
=\sum_{j=0}^nc_je^{-iy_jk\delta}\phi_X(y_j)$,}

\noindent dan dengan cara yang sama

\Centerline{$\Expn(g_k(Z))
=\sum_{j=0}^nc_je^{-iy_jk\delta}\phi_Z(y_j)$,}

\noindent sehingga

\Centerline{$|\Expn(g_k(X))-\Expn(g_k(Z))|
\le\sum_{j=0}^n|c_j||\phi_X(y_j)-\phi_Z(y_j)|
\le\sum_{j=0}^n|c_j|\eta\le\delta$.}

\noindent Selanjutnya,

\Centerline{$|h_k(x)-g_k(x)|\le\delta$ untuk setiap
$x\in[-b-(m+1)\delta+k\delta,b+m\delta+k\delta]\supseteq[-b,b]$,}

\Centerline{$|h_k(x)-g_k(x)|\le 2$ untuk setiap $x$,}

\Centerline{$\Pr(|X|\ge b)\le\Bover{M}{b}=\delta$,}

\noindent sehingga $\Expn(|h_k(X)-g_k(X)|)\le 3\delta$; dan demikian pula
$\Expn(|h_k(Z)-g_k(Z)|)\le 3\delta$. Dengan menggabungkan keduanya,

\Centerline{$|\Expn(h_k(X))-\Expn(h_k(Z))|\le 7\delta=\epsilon$}

\noindent bila $-m\le k\le m+1$.

\medskip

{\bf (c)} Sekarang andaikan $-b\le a\le b$. Maka ada $k$ sedemikian sehingga
$-m\le k\le m+1$ dan $a\le k\delta\le a+\delta$. Karena

\Centerline{$\chi\ocint{-\infty,a}
\le\chi\ocint{-\infty,k\delta}
\le h_k
\le\chi\ocint{-\infty,(k+1)\delta}
\le\chi\ocint{-\infty,a+2\delta}$,}

\noindent kita harus mempunyai

\Centerline{$\Pr(X\le a)\le\Expn(h_k(X))$,}

\Centerline{$\Expn(h_{k}(Z))\le\Pr(Z\le a+2\delta)\le\Pr(Z\le
a+\epsilon)$.}

\noindent Tetapi ini berarti bahwa

\Centerline{$\Pr(X\le a)\le\Expn(h_k(X))\le\Expn(h_k(Z))+\epsilon
\le\Pr(Z\le a+\epsilon)+\epsilon$}

\noindent bila $a\in[-b,b]$.

\medskip

{\bf (d)} Untuk kasus $a\ge b$, $a\le -b$, kita tentu mempunyai

\Centerline{$b(1-F_Z(b))=b\Pr(Z>b)\le\Expn(|Z|)\le M$,}

\noindent jadi jika $a\ge b$ maka

\Centerline{$F_X(a)\le 1\le F_Z(a)+1-F_Z(b)\le F_Z(a)+\Bover{M}{b}
=F_Z(a)+\delta\le F_Z(a+\epsilon)+\epsilon$.}

\noindent Demikian juga,

\Centerline{$bF_X(-b)\le \Expn(|X|)\le M$,}

\noindent Jadi

\Centerline{$F_X(a)\le\delta\le F_Z(a+\epsilon)+\epsilon$}

\noindent untuk setiap $a\le -b$. Ini menyelesaikan pembuktian.
}%end of proof of 285P

\vleader{72pt}{285Q}{Hukum kejadian langka: Teorema} Untuk setiap $M\ge 0$
dan $\epsilon>0$, ada $\delta>0$ sedemikian sehingga bila
$X_0,\ldots,X_n$ variabel-variabel acak bebas bernilai $\{0,1\}$ dengan
$\Pr(X_k=1)=p_k\le\delta$ untuk setiap $k\le n$ dan
$\sum_{k=0}^np_k=\lambda\le M$, serta $X=X_0+\ldots+X_n$, maka

\Centerline{$|\Pr(X=m)-\Bover{\lambda^m}{m!}e^{-\lambda}|\le\epsilon$}

\noindent untuk setiap $m\in\Bbb N$.

\proof{{\bf (a)} Kita mulai dengan menghitung beberapa fungsi karakteristik.
Pertama, fungsi karakteristik $\phi_k$ dari $X_k$ diberikan oleh

\Centerline{$\phi_k(y)=(1-p_k)e^{iy0}+p_ke^{iy1}=1+p_k(e^{iy}-1)$.}

\noindent Selanjutnya, jika $Z$ suatu variabel acak Poisson dengan parameter
$\lambda$ (yakni jika $\Pr(Z=m)=\lambda^me^{-\lambda}/m!$ untuk setiap
$m\in\Bbb N$; satu-satunya hal yang perlu diketahui saat ini mengenai
distribusi Poisson ialah bahwa
$\sum_{m=0}^{\infty}\lambda^me^{-\lambda}/m!=1$), maka fungsi
karakteristiknya $\phi_Z$ diberikan oleh

\Centerline{$\phi_Z(y)
=\sum_{m=0}^{\infty}\Bover{\lambda^m}{m!}e^{-\lambda}e^{iym}
=e^{-\lambda}\sum_{m=0}^{\infty}\Bover{(\lambda e^{iy})^m}{m!}
=e^{-\lambda}e^{\lambda e^{iy}}
=e^{\lambda(e^{iy}-1)}$.}

\medskip

{\bf (b)} Sebelum mengurus nilai $\delta$ dan $\eta$, saya tunjukkan cara
menaksir $\phi_X(y)-\phi_Z(y)$. Kita tahu bahwa

\Centerline{$\phi_X(y)=\prod_{k=0}^n\phi_k(y)$}

\noindent (menggunakan 285I), sementara

\Centerline{$\phi_Z(y)=\prod_{k=0}^ne^{p_k(e^{iy}-1)}$.}

\noindent Karena $\phi_k(y)$ dan $e^{p_k(e^{iy}-1)}$ semuanya bermodulus
paling besar $1$ (memang,

\Centerline{$|e^{p_k(e^{iy}-1)}|=e^{-p_k(1-\cos y)}\le 1$,)}

\noindent 285O mengatakan kepada kita bahwa

\Centerline{$|\phi_X(y)-\phi_Z(y)|
\le\sum_{k=0}^n|\phi_k(y)-e^{p_k(e^{iy}-1)}|
=\sum_{k=0}^n|e^{p_k(e^{iy}-1)}-1-p_k(e^{iy}-1)|$.}

\medskip

{\bf (c)} Jadi kita perlu melakukan sedikit analisis. Untuk menaksir
$|e^z-1-z|$ ketika $\Real z\le 0$, tinjau fungsi

\Centerline{$g(t)=\Real(c(e^{tz}-1-tz))$}

\noindent dengan $|c|=1$. Kita mempunyai $g(0)=g'(0)=0$ dan

\Centerline{$|g''(t)|=|\Real(c(z^2e^{tz}))|\le|c||z^2||e^{tz}|\le|z|^2$}

\noindent untuk setiap $t\ge 0$, sehingga

\Centerline{$|g(1)|\le\Bover12|z|^2$}

\noindent menurut teorema Taylor bernilai real dengan suku sisa, atau dengan
cara lain. Karena $c$ sembarang,

\Centerline{$|e^{z}-1-z|\le\Bover12|z|^2$}

\noindent bila $\Real z\le 0$. Khususnya,

\Centerline{$|e^{p_k(e^{iy}-1)}-1-p_k(e^{iy}-1)|
\le\Bover12p_k^2|e^{iy}-1|^2\le 2p_k^2$}

\noindent untuk setiap $k$, dan

\Centerline{$|\phi_X(y)-\phi_Z(y)|
\le\sum_{k=0}^n|e^{p_k(e^{iy}-1)}-1-p_k(e^{iy}-1)|
\le 2\sum_{k=0}^np_k^2$}

\noindent untuk setiap $y\in\Bbb R$.

\medskip

{\bf (d)} Sekarang untuk taksiran terperinci. Untuk $M\ge 0$ dan
$\epsilon>0$, pilih $\eta>0$ dan $y_0,\ldots,y_l\in\Bbb R$ sedemikian sehingga

\Centerline{$\Pr(X\le a)\le\Pr(Z\le a+\Bover12)+\Bover{\epsilon}2$}

\noindent bila $X$, $Z$ variabel-variabel acak bernilai real,
$\Expn(|X|)\le M$, $\Expn(|Z|)\le M$ dan
$|\phi_X(y_j)-\phi_Z(y_j)|\le\eta$ untuk setiap $j\le l$ (285P). Ambil
$\delta=\bover{\eta}{2M+1}$
dan andaikan $X_0,\ldots,X_n$ variabel-variabel acak bebas bernilai
$\{0,1\}$ dengan $\Pr(X_k=1)=p_k\le\delta$ untuk setiap $k\le n$, dengan
$\lambda=\sum_{k=0}^np_k$ tidak lebih dari $M$. Tetapkan
$X=X_0+\ldots+X_n$
dan misalkan $Z$ variabel acak Poisson dengan parameter $\lambda$. Maka,
menurut argumen (a)--(c),

\Centerline{$|\phi_X(y)-\phi_Z(y)|\le 2\sum_{k=0}^np_k^2
\le 2\delta\sum_{k=0}^np_k=2\delta\lambda\le\eta$}

\noindent untuk setiap $y\in\Bbb R$. Selain itu,

\Centerline{$\Expn(|X|)=\Expn(X)=\sum_{k=0}^np_k=\lambda\le M$,}

\Centerline{$\Expn(|Z|)=\Expn(Z)
=\sum_{m=0}^{\infty}m\Bover{\lambda^m}{m!}e^{-\lambda}
=e^{-\lambda}\sum_{m=1}^{\infty}\Bover{\lambda^m}{(m-1)!}
=e^{-\lambda}\sum_{m=0}^{\infty}\Bover{\lambda^{m+1}}{m!}
=\lambda\le M$.}

\noindent Jadi

\Centerline{$\Pr(X\le a)\le\Pr(Z\le a+\Bover12)+\Bover{\epsilon}2$,}

\Centerline{$\Pr(Z\le a)\le\Pr(X\le a+\Bover12)+\Bover{\epsilon}2$}

\noindent untuk setiap $a$. Tetapi karena $X$ dan $Z$ hanya mengambil nilai
dalam $\Bbb N$,

\Centerline{$|\Pr(X\le m)-\Pr(Z\le m)|\le\Bover{\epsilon}2$}

\noindent untuk setiap $m\in\Bbb N$, dan

\Centerline{$|\Pr(X=m)-\Bover{\lambda^m}{m!}e^{-\lambda}|
=|\Pr(X=m)-\Pr(Z=m)|
\le\epsilon$}

\noindent untuk setiap $m\in\Bbb N$, sesuai kebutuhan.
}%end of proof of 285Q

\leader{285R}{Konvolusi}\dvro{ Jika}{ Ingat dari 257A bahwa jika}
$\nu$, $\tilde\nu$ ukuran-ukuran probabilitas Radon pada $\BbbR^r$, maka
%\tilde\nu  not  \nuprime
%because we have  \tilde f  for its R-N deriviative
\dvro{$\phi_{\nu*\tilde\nu}(y)
=\phi_{\nu}(y)\phi_{\tilde\nu}(y)$
untuk setiap $y\in\BbbR^r$.}{keduanya mempunyai konvolusi
$\nu*\tilde\nu$ yang didefinisikan dengan

\Centerline{$(\nu*\tilde\nu)(E)=(\nu\times\tilde\nu)\{(x,y):x+y\in E\}$}

\noindent untuk setiap himpunan Borel $E\subseteq\BbbR^r$; konvolusi ini juga
merupakan ukuran probabilitas Radon. Menurut 257B, fungsi karakteristik
$\phi_{\nu*\tilde\nu}$ dapat dihitung langsung: kita mempunyai

$$\eqalign{\phi_{\nu*\tilde\nu}(y)
&=\int e^{iy\dotproduct x}(\nu*\tilde\nu)(dx)
=\int e^{iy\dotproduct (x+x')}\nu(dx)\tilde\nu(dx')\cr
&=\int e^{iy\dotproduct x}e^{iy\dotproduct x'}\nu(dx)\tilde\nu(dx')
=\int e^{iy\dotproduct x}\nu(dx)
  \int e^{iy\dotproduct x'}\tilde\nu(dx')
=\phi_{\nu}(y)\phi_{\tilde\nu}(y)}$$

\noindent untuk setiap $y\in\BbbR^r$. (Jadi konvolusi ukuran bersesuaian
dengan perkalian titik demi titik fungsi karakteristik, sebagaimana konvolusi
fungsi bersesuaian dengan perkalian titik demi titik transformasi Fourier.)
Karena jumlah variabel acak bebas bersesuaian dengan konvolusi distribusinya
(272T), ini memberi cara lain untuk memandang 285I. Ingat pula bahwa jika
$\nu$, $\tilde\nu$ mempunyai turunan Radon--Nikod\'ym $f$, $\tilde f$ terhadap
ukuran Lebesgue, maka $f*\tilde f$ merupakan turunan Radon--Nikod\'ym dari
$\nu*\tilde\nu$ (257F).
}%end of dvro

\leader{285S}{Topologi samar dan konvergensi titik demi titik fungsi
karakteristik}\dvro{ Tuliskan}{ Dalam 285L kita melihat bahwa barisan
$\sequencen{\nu_n}$ ukuran probabilitas Radon pada $\BbbR^r$
konvergen dalam topologi samar ke suatu ukuran probabilitas Radon $\nu$ jika
dan hanya jika

\Centerline{$\lim_{n\to\infty}\int e^{iy\dotproduct
x}\nu_n(dx)=\int
e^{iy\dotproduct x}\nu(dx)$}

\noindent untuk setiap $y\in\BbbR^r$; yakni, jika dan hanya jika

\Centerline{$\lim_{n\to\infty}\rho'_y(\nu_n,\nu)=0$ untuk setiap $y\in\Bbb
R^r$,}

\noindent dengan menulis}

\Centerline{$\rho'_y(\nu,\nuprime)=|\int e^{iy\dotproduct x}\nu(dx)
-\int e^{iy\dotproduct x}\nuprime(dx)|$}

\noindent untuk ukuran-ukuran probabilitas Radon $\nu$, $\nuprime$ pada
$\BbbR^r$ dan $y\in\BbbR^r$. \dvro{Tuliskan}{Wajar untuk menanyakan apakah
pseudometrik-pseudometrik $\rho'_y$ benar-benar mendefinisikan topologi samar.
Tuliskan} $\frak T$ untuk topologi samar dan $\frak S$ untuk topologi yang
didefinisikan oleh $\{\rho'_y:y\in\BbbR^r\}$\cmmnt{. Kita tentu mempunyai
$\frak S\subseteq\frak T$, karena setiap $\rho'_y$ merupakan salah satu
pseudometrik yang digunakan dalam definisi $\frak T$. Kita juga tahu bahwa
$\frak S$ dan $\frak T$ mempunyai barisan konvergen yang sama, serta bahwa
$\frak T$ termetrikkan (lihat 285Xt). Namun, semua ini belum berarti bahwa
kedua topologi tersebut sama; hasil berikut menunjukkan bahwa keduanya memang
berbeda.}%end of comment

\leader{285T}{Proposisi} Misalkan $y_0,\ldots,y_n\in\Bbb R$
dan $\eta>0$. Maka ada tak berhingga banyak $m\in\Bbb N$ sedemikian sehingga
$|1-e^{iy_km}|\le\eta$ untuk setiap $k\le n$.

\proof{ Pilih $\eta_1,\ldots,\eta_r\in\Bbb R$ sedemikian sehingga
$1=\eta_0,\eta_1,\ldots,\eta_r$ bebas linear atas $\Bbb Q$ dan setiap
$y_k/2\pi$ merupakan kombinasi linear dari $\eta_j$ atas $\Bbb Q$; tuliskan
$y_k=2\pi\sum_{j=0}^rq_{kj}\eta_j$ dengan setiap $q_{kj}\in\Bbb Q$.
Nyatakan $q_{kj}$ sebagai $p_{kj}/p$, dengan setiap $p_{kj}\in\Bbb Z$
dan $p\in\Bbb N\setminus\{0\}$. Tetapkan
$M=\max_{k\le n}\sum_{j=0}^r|p_{kj}|$.

Ambil sembarang $m_0\in\Bbb N$ dan pilih $\delta>0$ sedemikian sehingga
$|1-e^{2\pi ix}|\le\eta$ bila $|x|\le2\pi M\delta$. Menurut Teorema
Ekuidistribusi Weyl (281N), ada tak berhingga banyak $m$ sedemikian sehingga
$\fraction{m\eta_j}\le\delta$ bila $1\le j\le r$; khususnya, ada yang memenuhi
$m\ge m_0$. Tetapkan $m_j=\lfloor m\eta_j\rfloor$, sehingga
$|m\eta_j-m_j|\le\delta$ untuk $0\le j\le r$. Maka

\Centerline{$|mpy_k-2\pi\sum_{j=0}^rp_{kj}m_j|
\le2\pi\sum_{j=0}^r|p_{kj}||m\eta_j-m_j|
\le 2\pi M\delta$,}

\noindent jadi bahwa

\Centerline{$|1-e^{iy_kmp}|
=|1-\exp(i(mpy_k-2\pi\sum_{j=0}^rp_{kj}m_j))|
\le\eta$}

\noindent untuk setiap $k\le n$. Karena $mp\ge m_0$ dan $m_0$ sembarang,
hasilnya terbukti.
}%end of proof of 285T

\leader{285U}{Akibat} Topologi $\frak S$ dan $\frak T$ pada ruang ukuran
probabilitas Radon pada $\Bbb R$, sebagaimana dijelaskan dalam 285S, berbeda.

\proof{ Misalkan $\delta_x$ ukuran Dirac pada $\Bbb R$ yang terpusat di
$x$. Menurut 285T, setiap anggota $\frak S$ yang memuat $\delta_0$ juga memuat
$\delta_m$ untuk tak berhingga banyak $m\in\Bbb N$. Di sisi lain, himpunan

\Centerline{$G=\{\nu:\int e^{-x^2}\nu(dx)>\Bover12\}$}

\noindent merupakan anggota $\frak T$ yang memuat $\delta_0$, tetapi tidak
memuat $\delta_m$ untuk setiap bilangan bulat $m\ne 0$. Jadi
$G\in\frak T\setminus\frak S$ dan $\frak T\ne\frak S$.
}%end of proof of 285U


\exercises{
\leader{285X}{Latihan dasar (a)}
%\spheader 285Xa
Misalkan $\nu$ suatu ukuran probabilitas Radon pada $\BbbR^r$, dengan
$r\ge 1$, dan andaikan $\int\|x\|\nu(dx)<\infty$. Tunjukkan bahwa fungsi
karakteristik $\phi$ dari $\nu$ terdiferensialkan (dalam arti penuh 262Fa) dan
bahwa $\pd{\phi}{\eta_j}(y)=i\int\xi_je^{iy\dotproduct x}\nu(dx)$ untuk setiap
$j\le r$ dan $y\in\BbbR^r$, dengan $\xi_j$, $\eta_j$ menyatakan koordinat
$x$ dan $y$ seperti biasa.
%285A

\sqheader 285Xb Misalkan $\pmb{X}=(X_1,\ldots,X_r)$ suatu keluarga variabel
acak bernilai real, dengan fungsi karakteristik $\phi_{\pmb{X}}$. Tunjukkan
bahwa fungsi karakteristik $\phi_{X_j}$ dari $X_j$ diberikan oleh

\Centerline{$\phi_{X_j}(y)=\phi_{\pmb{X}}(ye_j)$ untuk setiap $y\in\Bbb
R$,}

\noindent dengan $e_j$ vektor satuan ke-$j$ dari $\BbbR^r$.
%285A

\sqheader 285Xc Misalkan $X$ variabel acak bernilai real dan $\phi_X$ fungsi
karakteristiknya. Tunjukkan bahwa

\Centerline{$\phi_{aX+b}(y)=e^{iyb}\phi_X(ay)$}

\noindent untuk setiap $a$, $b$, $y\in\Bbb R$.
%285A

\spheader 285Xd\dvAnew{2014}
Misalkan $X$ variabel acak bernilai real yang tidak konstan secara esensial,
dan $\phi$ fungsi karakteristiknya. Tunjukkan bahwa $|\phi(y)|<1$ untuk semua
$y\in\Bbb R$, kecuali mungkin sejumlah terhitung banyaknya.
\Hint{Dukungan (256Xf) distribusi $X$ memuat dua titik berbeda $x$, $x'$;
jika $e^{iyx}\ne e^{iyx'}$, maka $|\phi(y)|<1$.}
%285D

\spheader 285Xe Misalkan $X$ variabel acak bernilai real dan
$\phi$ fungsi karakteristiknya.

\quad{(i)} Tunjukkan bahwa untuk setiap fungsi terintegralkan bernilai kompleks
$h$ pada $\Bbb R$,

\Centerline{$\Expn(\varhat{h}(X))=\Bover1{\sqrt{2\pi}}
\int_{-\infty}^{\infty}\phi(-y)h(y)dy$,}

\noindent dengan $\varhat{h}$ menyatakan transformasi Fourier dari $h$.

\quad{(ii)} Tunjukkan bahwa untuk setiap fungsi uji yang menurun cepat $h$,

\Centerline{$\Expn(h(X))
=\Bover1{\sqrt{2\pi}}\int_{-\infty}^{\infty}
  \phi(y)\varhat{h}(y)dy$.}
%285D

\spheader 285Xf Misalkan $\nu$ suatu ukuran probabilitas Radon pada
$\Bbb R$, dan andaikan fungsi karakteristiknya $\phi$ terintegralkan-kuadrat.
Tunjukkan bahwa $\nu$ merupakan ukuran integral tak tentu terhadap ukuran
Lebesgue dan turunan Radon--Nikod\'ym-nya juga terintegralkan-kuadrat.
\Hint{Gunakan 284O untuk mencari fungsi terintegralkan-kuadrat $f$ sedemikian
sehingga $\int f\times h=\bover1{\sqrt{2\pi}}\int\phi\times\varhat{h}$ untuk
setiap fungsi uji yang menurun cepat $h$, serta gagasan dari pembuktian 284G
untuk menunjukkan bahwa $\int_a^bf=\nu\ooint{a,b}$ bila $a<b$ dalam
$\Bbb R$.}
%285D 285Xe

\spheader 285Xg\dvAformerly{2{}85Yp} Misalkan $\nu$ suatu ukuran probabilitas
Radon pada $\BbbR^r$ dengan dukungan terbatas (definisi: 256Xf). Tunjukkan bahwa
fungsi karakteristiknya mulus.
%285D

\spheader 285Xh Misalkan $X$ suatu variabel acak normal dengan ekspektasi $a$
dan varians $\sigma^2$. Tunjukkan bahwa
$\Expn(e^X)=\exp(a+\bover12\sigma^2)$.
%285E

\sqheader 285Xi Misalkan $\pmb{X}=(X_1,\ldots,X_r)$ suatu keluarga variabel
acak bernilai real dengan fungsi karakteristik $\phi_{\pmb{X}}$. Andaikan
$\phi_{\pmb{X}}$ dapat dinyatakan dalam bentuk

\Centerline{$\phi_{\pmb{X}}(y)=\prod_{j=1}^r\phi_j(\eta_j)$}

\noindent untuk fungsi-fungsi $\phi_1,\ldots,\phi_r$, dengan menulis
$y=(\eta_1,\ldots,\eta_r)$ seperti biasa. Tunjukkan bahwa $X_1,\ldots,X_r$
saling bebas. \Hint{Tunjukkan bahwa $\phi_j$ harus merupakan kelipatan fungsi
karakteristik $X_j$; kemudian tunjukkan bahwa distribusi $\pmb{X}$ mempunyai
fungsi karakteristik yang sama dengan produk distribusi $X_j$.}
%285I

\spheader 285Xj Misalkan $X_1$, $X_2$ variabel acak bernilai real yang bebas
dengan distribusi yang sama, dan $\phi$ fungsi karakteristik $X_1-X_2$.
Tunjukkan bahwa $\phi(t)=\phi(-t)\ge 0$ untuk setiap
$t\in\Bbb R$.
%285I

\spheader 285Xk Misalkan $\nu$ suatu ukuran probabilitas Radon pada
$\Bbb R$, dengan fungsi karakteristik $\phi$. Tunjukkan bahwa

\Centerline{$\Bover12(\nu[c,d]+\nu\ooint{c,d})
=\Bover{i}{2\pi}\lim_{a\to\infty}
  \int_{-a}^a\Bover{e^{-idy}-e^{-icy}}{y}\phi(y)dy$}

\noindent bila $c<d$ dalam $\Bbb R$.
\Hint{Gunakan bagian (a) dari pembuktian 283F.}
%285J

\spheader 285Xl Misalkan $X$ variabel acak bernilai real dan
$\phi_X$ fungsi karakteristiknya. Tunjukkan bahwa

\Centerline{$\Pr(|X|\ge a)
\le 7a\int_0^{1/a}(1-\Real(\phi_X(y))dy$}

\noindent untuk setiap $a>0$.
%285J

\spheader 285Xm Kita mengatakan bahwa suatu himpunan $Q$ ukuran probabilitas
Radon pada $\Bbb R$ {\bf ketat secara seragam} jika untuk setiap $\epsilon>0$
ada $M\ge 0$ sedemikian sehingga
$\nu(\Bbb R\setminus[-M,M])\le\epsilon$ untuk setiap $\nu\in Q$. Tunjukkan
bahwa jika $Q$ sembarang keluarga ukuran probabilitas Radon pada $\Bbb R$ yang
ketat secara seragam, dan $\epsilon>0$, maka ada $\eta>0$ dan
$y_0,\ldots,y_n\in\Bbb R$ sedemikian sehingga

\Centerline{$\nu\ocint{-\infty,a}
\le\nuprime\ocint{-\infty,a+\epsilon}+\epsilon$}

\noindent bila $\nu$, $\nuprime\in Q$ dan
$|\phi_{\nu}(y_j)-\phi_{\nuprime}(y_j)|\le\eta$ untuk setiap $j\le n$, dengan
$\phi_{\nu}$ menyatakan fungsi karakteristik $\nu$.
%285M

\spheader 285Xn Misalkan $\sequencen{\nu_n}$ suatu barisan ukuran probabilitas
Radon pada $\Bbb R$, dan andaikan barisan ini konvergen dalam topologi samar ke
ukuran probabilitas Radon $\nu$. Tunjukkan bahwa
$\{\nu\}\cup\{\nu_n:n\in\Bbb N\}$ ketat secara seragam dalam arti
285Xm.
%285M

\sqheader 285Xo Misalkan $\nu$, $\nuprime$ dua ukuran Radon berhingga total
pada $\BbbR^r$ yang sama pada semua setengah-ruang tertutup, yakni himpunan
berbentuk $\{x:x\dotproduct y\ge c\}$ untuk $c\in\Bbb R$, $y\in\BbbR^r$.
Tunjukkan bahwa $\nu=\nuprime$. \Hint{Reduksikan pada kasus
$\nu\BbbR^r=\nuprime\Bbb R^r=1$ dan gunakan 285M.}
%285M

\sqheader 285Xp\dvAformerly{2{}85Xm}
Untuk $\gamma>0$, {\bf distribusi Cauchy} dengan pusat $0$ dan parameter skala
$\gamma$ ialah ukuran probabilitas Radon $\nu_{\gamma}$ yang didefinisikan oleh
rumus

\Centerline{$\nu_{\gamma}(E)
=\Bover{\gamma}{\pi}\int_E\Bover1{\gamma^2+t^2}dt$.}

\noindent (i) Tunjukkan bahwa jika $X$ variabel acak dengan distribusi
$\nu_{\gamma}$, maka $\Pr(X\ge 0)=\Pr(|X|\ge\gamma)=\bover12$. (ii) Tunjukkan
bahwa fungsi karakteristik $\nu_{\gamma}$ ialah
$y\mapsto e^{-\gamma|y|}$. \Hint{283Xq.} (iii) Tunjukkan bahwa jika $X$ dan
$Y$ variabel-variabel acak bebas dengan distribusi Cauchy, keduanya berpusat di
$0$ dan masing-masing mempunyai parameter skala $\gamma$, $\delta$, serta
$\alpha$, $\beta$ tidak keduanya $0$, maka $\alpha X+\beta Y$ mempunyai
distribusi Cauchy berpusat di $0$ dengan parameter skala
$|\alpha|\gamma+|\beta|\delta$. (iv) Tunjukkan bahwa jika $X$ dan $Y$ variabel
acak bebas berdistribusi normal dengan ekspektasi $0$, maka $X/Y$ mempunyai
distribusi Cauchy.
%285M

\sqheader 285Xq Misalkan $X_1,X_2,\ldots$ suatu barisan variabel acak bebas
dan berdistribusi identik, semuanya dengan ekspektasi nol dan varians $1$;
misalkan $\phi$ fungsi karakteristik bersamanya. Untuk setiap $n\ge 1$,
tetapkan $S_n=\bover1{\sqrt{n}}(X_1+\ldots+X_n)$.

\quad{(i)} Tunjukkan bahwa fungsi karakteristik $\phi_n$ dari $S_n$
diberikan dengan rumus $\phi_n(y)=(\phi(\Bover{y}{\sqrt{n}}))^{n}$
untuk setiap $n$.

\quad{(ii)} Tunjukkan bahwa
$|\phi_n(y)-e^{-y^2/2}|\le n|\phi(\Bover{y}{\sqrt{n}})-e^{-y^2/2n}|$
untuk $n\ge 1$ dan $y\in\Bbb R$.

\ifdim\pagewidth>467pt\fontdimen3\tenrm=3.33pt\fi
\quad(iii) Dengan menetapkan $h(y)=\phi(y)-e^{-y^2/2}$, tunjukkan bahwa
$h(0)=h'(0)=h''(0)=0$, dan karena itu
$\lim_{n\to\infty} nh(y/\sqrt{n})=0$, sehingga
$\lim_{n\to\infty}\phi_n(y)=e^{-y^2/2}$ untuk setiap $y\in\Bbb R$.
\fontdimen3\tenrm=1.67pt

\quad{(iv)} Tunjukkan bahwa $\lim_{n\to\infty}\Pr(S_n\le a)
=\Bover1{\sqrt{2\pi}}\int_{-\infty}^ae^{-x^2/2}dx$ untuk setiap
$a\in\Bbb R$.
%285N

\sqheader 285Xr\dvAformerly{2{}85Xm}
Variabel acak $X$ mempunyai {\bf distribusi Poisson}
dengan parameter $\lambda>0$ jika $\Pr(X=n)=e^{-\lambda}\lambda^n/n!$ untuk
setiap $n\in\Bbb N$. (i) Tunjukkan bahwa dalam hal ini
$\Expn(X)=\Var(X)=\lambda$. (ii) Tunjukkan bahwa jika $X$ dan $Y$ variabel
acak bebas dengan distribusi Poisson, maka $X+Y$ mempunyai distribusi Poisson.
(iii) Carilah pembuktian (ii) berdasarkan 285Q.
%285Q

\sqheader 285Xs Untuk $x\in\BbbR^r$, misalkan $\delta_x$ ukuran Dirac pada
$\BbbR^r$ yang terpusat di $x$. Tunjukkan bahwa
$\delta_x*\delta_y=\delta_{x+y}$ untuk semua $x$, $y\in\BbbR^r$.
%285R

\spheader 285Xt Misalkan $P$ himpunan ukuran probabilitas Radon pada
$\BbbR^r$. Untuk $y\in\BbbR^r$, tetapkan
$\rho'_y(\nu,\nuprime)=|\phi_{\nu}(y)-\phi_{\nuprime}(y)|$ untuk semua $\nu$,
$\nuprime\in P$, menulis $\phi_{\nu}$ untuk fungsi karakteristik dari
$\nu$. Tetapkan $\psi(x)=\bover1{(\sqrt{2\pi})^r}e^{-x\dotproduct x/2}$ untuk
$x\in\BbbR^r$. Tunjukkan bahwa topologi samar pada $P$ didefinisikan oleh
keluarga $\{\rho_{\psi}\}\cup\{\rho'_y:y\in\Bbb Q^r\}$, dengan
$\rho_{\psi}$ didefinisikan seperti dalam 285K, dan karenanya termetrikkan.
\Hint{281K; bandingkan 285Xm.}
%285S

\sqheader 285Xu\dvAformerly{2{}85Xr}
Misalkan $\phi:\BbbR^r\to\Bbb C$ fungsi karakteristik suatu ukuran
probabilitas Radon pada $\BbbR^r$. Tunjukkan bahwa $\phi(0)=1$ dan
$\sum_{j=0}^n\sum_{k=0}^nc_j\bar c_k\phi(a_j-a_k)\ge 0$ bila
$a_0,\ldots,a_n\in\BbbR^r$ dan $c_0,\ldots,c_n\in\Bbb C$. (`Teorema
Bochner' menyatakan bahwa syarat-syarat ini tidak hanya perlu, tetapi juga cukup,
agar $\phi$ menjadi fungsi karakteristik; lihat 445N dalam Jilid 4.)
%285*

\leader{285Y}{Latihan lebih lanjut (a)}
%\spheader 285Ya
Misalkan $\nu$ suatu ukuran probabilitas Radon pada $\BbbR^r$. Tuliskan

\Centerline{$\varhat{\nu}(y)=\Bover1{(\sqrt{2\pi})^r}
\int e^{-iy\dotproduct x}\nu(dx)$}

\noindent untuk setiap $y\in\BbbR^r$.

\quad{(i)} Dengan $\phi_{\nu}$ menyatakan fungsi karakteristik dari
$\nu$, tunjukkan bahwa $\varhat{\nu}(y)
=\Bover1{(\sqrt{2\pi})^r}\phi_{\nu}(-y)$ untuk setiap $y\in\BbbR^r$.

\quad{(ii)} Tunjukkan bahwa
$\int h(y)\varhat{\nu}(y)dy=\int\varhat h(x)\nu(dx)$
untuk setiap fungsi Lebesgue terintegralkan bernilai kompleks $h$ pada
$\BbbR^r$, dengan transformasi Fourier $\varhat{h}$ didefinisikan seperti
dalam 283Wa.

\quad{(iii)} Tunjukkan bahwa $\int h(x)\nu(dx)
=\int\varcheck{h}(y)\varhat{\nu}(y)dy$ untuk setiap fungsi uji yang menurun
cepat $h$ pada $\BbbR^r$.

\quad{(iv)} Tunjukkan bahwa jika $\nu$ merupakan ukuran integral tak tentu
terhadap ukuran Lebesgue, dengan turunan Radon--Nikod\'ym $f$, maka
$\varhat{\nu}$ merupakan transformasi Fourier dari $f$.
%285D

\spheader 285Yb Misalkan $\nu$ suatu ukuran probabilitas Radon pada $\BbbR^r$,
dengan fungsi karakteristik $\phi$. Tunjukkan bahwa bila $c\le d$ dalam
$\BbbR^r$, maka

$$\bigl(\bover{i}{2\pi}\bigr)^r\lim_{\alpha_1,\ldots,\alpha_r\to\infty}
\int_{[-a,a]}\bigl(\prod_{j=1}^r
\Bover{e^{-i\delta_j\eta_j}-e^{-i\gamma_j\eta_j}}{\eta_j}\bigr)
\phi(y)dy$$

\noindent ada dan nilainya terletak di antara $\nu\ooint{c,d}$ dan
$\nu[c,d]$, dengan menulis $a=(\alpha_1,\ldots,\alpha_r)$ dan
$\ooint{c,d}=\prod_{j\le r}\ooint{\gamma_j,\delta_j}$ jika
$c=(\gamma_1,\ldots,\gamma_r)$ dan $d=(\delta_1,\ldots,\delta_r)$.
%285Xk, 285J, 283Wc

\spheader 285Yc Misalkan $\sequencen{X_n}$ suatu barisan variabel acak bebas
dan berdistribusi identik (yang tidak konstan secara esensial). Tunjukkan bahwa
$\lim_{n\to\infty}\Pr(|\sum_{k=0}^nX_k|\le\alpha)=0$ untuk setiap
$\alpha\in\Bbb R$.
%285Xd, 285Xk, 285J

\spheader 285Yd Untuk ukuran-ukuran probabilitas Radon $\nu$, $\nuprime$ pada
$\BbbR^r$, tetapkan

$$\eqalign{\rho(\nu,\nuprime)=\inf\{\epsilon:\epsilon\ge 0,\,
\nu\ocint{-\infty,a}\le\nuprime\ocint{-\infty,a+\epsilon\tbf{1}}+\epsilon
&\le\nu\ocint{-\infty,a+2\epsilon\tbf{1}}+2\epsilon\cr
&\text{untuk setiap }a\in\BbbR^r\},\cr}$$

\noindent dengan menulis $\ocint{-\infty,a}
=\{(\xi_1,\ldots,\xi_r):\xi_j\le\alpha_j$ untuk setiap $j\le r\}$
bila $a=(\alpha_1,\ldots,\alpha_r)$, dan
$\tbf{1}=(1,\ldots,1)\in\BbbR^r$. Tunjukkan bahwa $\rho$ merupakan metrik
pada himpunan ukuran probabilitas Radon pada $\BbbR^r$, dan topologi yang
didefinisikannya ialah topologi samar. (Bandingkan 274Yc.)
%285K

\spheader 285Ye Misalkan $r\ge 1$ dan $P$ himpunan ukuran probabilitas Radon
pada $\BbbR^r$. Untuk $m\in\Bbb N$, misalkan $\rho^*_m$ pseudometrik pada
$P$ yang didefinisikan dengan
$\rho^*_m(\nu,\nuprime)=\sup_{\|y\|\le m}|\phi_{\nu}(y)-\phi_{\nuprime}(y)|$
untuk $\nu$, $\nuprime\in P$, dengan $\phi_{\nu}$ menyatakan fungsi
karakteristik $\nu$. Tunjukkan bahwa $\{\rho^*_m:m\in\Bbb N\}$ mendefinisikan
topologi samar pada $P$.
%285K

\spheader 285Yf Misalkan $r\ge 1$. Kita mengatakan bahwa suatu himpunan $Q$
ukuran probabilitas Radon pada $\BbbR^r$ {\bf ketat secara seragam} jika untuk
setiap $\epsilon>0$ ada himpunan kompak $K\subseteq\BbbR^r$ sedemikian sehingga
$\nu(\BbbR^r\setminus K)\le\epsilon$ untuk setiap $\nu\in Q$. Tunjukkan bahwa
jika $Q$ sembarang keluarga ukuran probabilitas Radon pada $\BbbR^r$ yang ketat
secara seragam, dan $\epsilon>0$, maka ada $\eta>0$,
$y_0,\ldots,y_n\in\BbbR^r$ sedemikian sehingga
$\nu\ocint{-\infty,a}\le\nuprime\ocint{-\infty,a+\epsilon\tbf{1}}+\epsilon$
bila $\nu$, $\nuprime\in Q$ dan $a\in\BbbR^r$, serta
$|\phi_{\nu}(y_j)-\phi_{\nuprime}(y_j)|\le\eta$ untuk setiap $j\le n$, dengan
$\phi_{\nu}$ menyatakan fungsi karakteristik $\nu$.
%285M, 285Xm

\spheader 285Yg Tunjukkan bahwa untuk setiap $M\ge 0$, himpunan ukuran
probabilitas Radon $\nu$ pada $\BbbR^r$ yang memenuhi
$\int\|x\|\nu(dx)\le M$ ketat secara seragam dalam arti 285Yf.
%285Yf, 285Xm, 285M

\spheader 285Yh Misalkan $C_b(\BbbR^r)$ ruang Banach fungsi-fungsi kontinu
terbatas bernilai real pada $\BbbR^r$.

\quad{(i)} Tunjukkan bahwa setiap ukuran probabilitas Radon $\nu$ pada
$\Bbb R^r$ bersesuaian dengan fungsional linear kontinu
$h_{\nu}:C_b(\Bbb R^r)\to\Bbb R$, dengan $h_{\nu}(f)=\int fd\nu$ untuk
$f\in C_b(\Bbb R^r)$.

\quad{(ii)} Tunjukkan bahwa jika $h_{\nu}=h_{\nuprime}$, maka $\nu=\nuprime$.

\quad{(iii)} Tunjukkan bahwa topologi samar pada himpunan ukuran probabilitas
Radon bersesuaian dengan topologi lemah* pada ruang dual
$(C_b(\BbbR^r))^*$ dari $C_b(\BbbR^r)$.
%285M mt28bits

\spheader 285Yi Misalkan $r\ge 1$ dan $P$ himpunan ukuran probabilitas Radon
pada $\BbbR^r$. Untuk $m\in\Bbb N$, misalkan $\tilde\rho^*_m$ pseudometrik
pada $P$ yang didefinisikan dengan

\Centerline{$\tilde\rho^*_m(\nu,\nuprime)
=\int_{\{y:\|y\|\le m\}}|\phi_{\nu}(y)-\phi_{\nuprime}(y)|dy$}

\noindent untuk $\nu$, $\nuprime\in P$, dengan $\phi_{\nu}$ menyatakan fungsi
karakteristik $\nu$. Tunjukkan bahwa $\{\tilde\rho^*_m:m\in\Bbb N\}$
mendefinisikan topologi samar pada $P$.
%285M mt28bits

\spheader 285Yj Misalkan $(\Omega,\Sigma,\mu)$ suatu ruang probabilitas.
Andaikan $\sequencen{X_n}$ barisan variabel acak bernilai real pada $\Omega$,
dan $X$ variabel acak bernilai real lainnya pada $\Omega$; misalkan
$\phi_{X_n}$, $\phi_X$ fungsi karakteristik yang bersesuaian. Tunjukkan bahwa
pernyataan-pernyataan berikut ekuivalen: (i)
$\lim_{n\to\infty}\Expn(h(X_n))=\Expn(h(X))$ untuk setiap fungsi kontinu
terbatas $h:\Bbb R\to\Bbb R$; (ii)
$\lim_{n\to\infty}\phi_{X_n}(y)=\phi_X(y)$ untuk setiap $y\in\Bbb R$.
%285Xm 285M

\spheader 285Yk Misalkan $(\Omega,\Sigma,\mu)$ suatu ruang probabilitas, dan
$P$ himpunan ukuran probabilitas Radon pada $\Bbb R$. (i) Tunjukkan bahwa ada
fungsi $\psi:L^0(\mu)\to P$ yang didefinisikan dengan menyatakan bahwa
$\psi(X^{\ssbullet})$ ialah distribusi $X$ bila $X$ variabel acak bernilai real
pada $\Omega$. (ii) Tunjukkan bahwa $\psi$ kontinu untuk topologi konvergensi
dalam ukuran pada $L^0(\mu)$ dan topologi samar pada $P$. (Bandingkan 271Yd.)
%285Yj 285M

\spheader 285Yl Misalkan $X$ variabel acak bernilai real dengan varians
berhingga. Tunjukkan bahwa untuk setiap $\eta\ge 0$,

\Centerline{$|\phi(y)-1-iy\Expn(X)+\Bover12y^2\Expn(X^2)|
\le\Bover16\eta|y|^3\Expn(X^2)
+y^2\Expn(\psi_{\eta}(X))$,}

\noindent dengan $\phi$ menyatakan fungsi karakteristik $X$ dan
$\psi_{\eta}(x)=0$ untuk $|x|\le\eta$, $x^2$ untuk $|x|>\eta$.
%285N

\spheader 285Ym Misalkan $\epsilon\ge\delta>0$ dan
$X_0,\ldots,X_n$ variabel-variabel acak bernilai real yang bebas sedemikian
sehingga

\Centerline{$\Expn(X_k)=0$ untuk setiap $k\le n$,
\quad$\sum_{k=0}^n\Var(X_k)=1$,
\quad$\sum_{k=0}^n\Expn(\psi_{\delta}(X_k))\le\delta$}

\noindent (dengan menulis $\psi_{\delta}(x)=0$ jika $|x|\le\delta$, dan $x^2$
jika $|x|>\delta$). Tetapkan $\gamma=\epsilon/\sqrt{\delta^2+\delta}$, dan
misalkan $Z$ suatu variabel acak normal standar. Tunjukkan bahwa

\Centerline{$|\phi(y)-e^{-y^2/2}|
\le\Bover13\epsilon|y|^3+y^2(\delta+\Expn(\psi_{\gamma}(Z)))$}

\noindent untuk setiap $y\in\Bbb R$, dengan $\phi$ menyatakan fungsi
karakteristik dari $X=\sum_{k=0}^nX_k$.
% mt28bits 285Yl, 285N

\spheader 285Yn Tunjukkan bahwa untuk setiap $\epsilon>0$ ada
$\delta>0$ sedemikian sehingga bila $X_0,\ldots,X_n$ variabel-variabel acak
bernilai real yang bebas dan memenuhi

\Centerline{$\Expn(X_k)=0$ untuk setiap $k\le n$,
\quad$\sum_{k=0}^n\Var(X_k)=1$,
\quad$\sum_{k=0}^n\Expn(\psi_{\delta}(X_k))\le\delta$}

\noindent (dengan menulis $\psi_{\delta}(x)=0$ jika $|x|\le\delta$, dan $x^2$
jika $|x|>\delta$), maka
$|\phi(y)-e^{-y^2/2}|\le\epsilon(y^2+|y^3|)$ untuk setiap $y\in\Bbb R$,
dengan $\phi$ menyatakan fungsi karakteristik dari $X=X_0+\ldots+X_n$.
%285Ym, 285N

\spheader 285Yo Gunakan 285Yn untuk membuktikan teorema Lindeberg (274F).
%285Yn, 285N

\spheader 285Yp Misalkan $r\ge 1$ dan $P$ himpunan ukuran probabilitas Radon
pada $\BbbR^r$. Tunjukkan bahwa konvolusi, dipandang sebagai pemetaan dari
$P\times P$ ke $P$, kontinu bila $P$ diberi topologi samar.
%285R \Hint{281Xa and 257B will help.}

\spheader 285Yq Misalkan $\frak S$ topologi pada $\Bbb R$ yang didefinisikan
oleh $\{\rho'_y:y\in\Bbb R\}$, dengan
$\rho'_y(x,x')=|e^{iyx}-e^{iyx'}|$ (bandingkan 285S). Tunjukkan bahwa
penjumlahan dan pengurangan kontinu untuk $\frak S$ dalam arti 2A5A.
%285S

\spheader 285Yr\dvAnew{2014} Misalkan $\nu$ suatu ukuran probabilitas pada
$\Bbb R$. Tunjukkan bahwa
$|\phi_{\nu}(y)-\phi_{\nu}(y')|^2\le 2(1-\Real\phi_{\nu}(y-y'))$ untuk setiap
$y$, $y'\in\Bbb R$.
%|\phi_{\nu}(y)-\phi_{\nu}(y')|^2\le\int|e^{iyx)-e^{iy'x}|^2\nu(dx)
%285J

\spheader 285Ys\dvAnew{2014} Misalkan $\sequencen{\nu_n}$ suatu barisan ukuran
probabilitas pada $\Bbb R$. Tetapkan $E=\{y:y\in\Bbb R$,
$\lim_{n\to\infty}\phi_{\nu_n}(y)=1\}$. (i) Tunjukkan bahwa $E-E$ dan
$E+E$ termuat dalam $E$. (ii) Tunjukkan bahwa jika $E$ tidak terabaikan
terhadap ukuran Lebesgue, maka himpunan itu adalah seluruh $\Bbb R$.
%285Yr 285J

\spheader 285Yt\dvAformerly{2{}85Xs}
Misalkan $\sequencen{X_n}$ suatu barisan variabel acak bernilai real yang saling
bebas, dan tetapkan $S_n=\sum_{j=0}^nX_j$ untuk setiap $n\in\Bbb N$.
Andaikan barisan distribusi $\sequencen{\nu_{S_n}}$ konvergen dalam topologi
samar ke suatu distribusi. Tunjukkan bahwa $\sequencen{S_n}$ konvergen dalam
ukuran, dan karena itu konvergen hampir di mana-mana.
%\Hint{285J, 273B, 285Ys.}
%285I  char fns of  X_n  converge
%  to  1  whenever char fns of  S_n  don't converge to  0
}%end of exercises

\endnotes{
\Notesheader{285} Seperti halnya transformasi Fourier, kekuatan metode yang
menggunakan fungsi karakteristik distribusi bertumpu pada tiga hal: (i) fungsi
karakteristik suatu distribusi menentukan distribusi itu (285M); (ii) sifat
distribusi yang kita minati tercermin pada sifat fungsi karakteristiknya yang
mudah dianalisis (285G, 285I, 285J); (iii) sifat-sifat fungsi karakteristik
tersebut benar-benar {\it berbeda} dari sifat-sifat yang bersesuaian pada
distribusi, sehingga dapat ditelaah dengan cara lain. Yang terpenting, fakta
bahwa (untuk barisan!) konvergensi distribusi dalam topologi samar bersesuaian
dengan konvergensi titik demi titik fungsi karakteristik (285L) memberi kita
jalan menuju teorema limit klasik, seperti dalam 285Q dan 285Xq. Dalam
285S-285U saya menunjukkan bahwa hasil untuk barisan ini tidak langsung
bersesuaian dengan karakterisasi alternatif topologi samar, meskipun hasil itu
dapat diadaptasi dengan lebih dari satu cara untuk memberikan karakterisasi
semacam itu (lihat 285Ye-285Yi).

Mengenai Teorema Limit Pusat, ada satu perbedaan mencolok antara metode yang
disarankan di sini dan metode dalam \S274. Pendekatan sebelumnya setidaknya
menawarkan kemungkinan teoretis untuk memberikan rumus eksplisit bagi
$\delta$ dalam 274F sebagai fungsi dari $\epsilon$, dan dengan demikian
taksiran laju konvergensi yang dapat diharapkan dalam Teorema Limit Pusat.
Argumen dalam bab ini melibatkan argumen kekompakan yang sepenuhnya
nonkonstruktif dalam 281A, sehingga tidak memberi kita cara untuk memperoleh
taksiran semacam itu. Namun, sebenarnya metode fungsi karakteristik, setelah
disempurnakan secara tepat, merupakan dasar taksiran terbaik yang diketahui,
seperti dalam teorema Berry--Ess\'een (274Hc).

Dalam 285D saya mencoba menunjukkan bagaimana fungsi karakteristik
$\phi_{\nu}$ dari suatu ukuran probabilitas Radon dapat dihubungkan dengan
`transformasi Fourier' $\varhat{\nu}$ dari $\nu$ yang langsung bersesuaian
dengan transformasi Fourier fungsi yang dibahas dalam \S\S283-284. Jika $f$
fungsi Lebesgue terintegralkan tak negatif dan kita mengambil $\nu$ sebagai
ukuran integral tak tentu yang bersesuaian, maka
$\varhat{\nu}=\varhatf$. Jadi konsep `transformasi Fourier suatu ukuran'
merupakan perluasan alami transformasi Fourier fungsi terintegralkan. Dilihat
dari arah sebaliknya, rumus dalam 285Dc menunjukkan bahwa $\nu$ dapat dipandang
sebagai representasi transformasi Fourier invers dari $\varhat{\nu}$ dalam
arti 284H-284I. Dengan mengambil $\nu$ sebagai ukuran yang memberi massa $1$
pada titik $0$, kita memperoleh fungsi delta Dirac, dengan transformasi Fourier
berupa fungsi konstan $\chi\Bbb R$. Gagasan-gagasan ini dapat diperluas tanpa
kesulitan untuk menangani konvolusi ukuran (285R).

Menarik bahwa, meskipun tidak ada karakterisasi memuaskan bagi fungsi-fungsi
yang merupakan transformasi Fourier dari fungsi terintegralkan, ada
karakterisasi bagi fungsi karakteristik distribusi probabilitas. Inilah
`teorema Bochner'. Saya memberikan syaratnya dalam 285Xu dan meminta Anda
membuktikan keperluannya sebagai latihan; kita sudah mempunyai tiga perempat
perangkat untuk membuktikan kecukupannya, tetapi langkah terakhir harus
menunggu Jilid 4.
}

\discrpage
